
% SGA 3，报告 VII_B；英文重建本第一轮的中文译本。
% 范围： 节 4 标题和开头约定，以及 节 4.1
% 全文, 依据文献局部页 73--74; 止于下节之前： 节 4.1.1.
% 合并阅读本页： 586--587. 重版印刷页： 575--576.
% 范围内原版边码： 541.
% 依据： Polo--Gille Exp7B-13oct24.pdf,
% SHA-256 D303690E0195654BBE7C3DAA5D89EA23E24CD709A3575217234AC19AE8836694.
% 参照了 Jacob Reinhold 的 CC BY 4.0 Markdown，作为对照来源。
% 本组件无图；所有数学内容均为原生 LaTeX。
\section*{4. 特征 \(p>0\) 特有的现象}
\phantomsection
\label{sga3:VIIB:4}
\addcontentsline{toc}{section}{4. 特征 p > 0 特有的现象}
在\SGARefLink{sga3:VIIB:4}{第 4 节}中，\(p\) 始终表示素数，\(k\) 表示特征 \(p\) 的伪紧环，且
\[
  \pi:k\longrightarrow k,\qquad x\longmapsto x^p,
\]
表示 \(k\) 的连续 Frobenius 自同态。
\subsection*{4.1}
\phantomsection
\label{sga3:VIIB:4.1}
设 \(X\) 是仿射代数为 \(A\) 的形式 \(k\)-簇。用 \(X^{(p/k)}\)，或简记 \(X^{(p)}\)，表示通过基变换 \(\pi:k\to k\)\textup{（参见\ \SGARefLink{sga3:VIIB:1.2.D}{1.2.D}）}得到的形式 \(k\)-簇
\[
  X\widehat\otimes_{\pi}k
\]
。其仿射代数是完备化张量积
\[
  A\widehat\otimes_{\pi}k.
\]
因此，连续态射
\[
  A\widehat\otimes_{\pi}k\longrightarrow A,
  \qquad
  a\widehat\otimes_{\pi}\lambda
  \longmapsto
  a^p\lambda,
\]
诱导形式 \(k\)-簇态射
\[
  \operatorname{Fr}(X/k):
  X\longrightarrow X^{(p/k)}.
\]
我们称 \(\operatorname{Fr}(X/k)\) 为 \emph{\(X\) 相对于 \(k\) 的 Frobenius 态射}；以下经常将 \(\operatorname{Fr}(X/k)\) 简记为 \(\operatorname{Fr}\)。
